\documentclass{beamer}
\usetheme{Madrid}
\usecolortheme{default}
\usepackage{fontspec}
\usepackage{polyglossia}
\setdefaultlanguage{english}
\usepackage{tikz}
\usepackage{hyperref}
\usepackage{graphicx}
\usepackage{booktabs}

% Custom colors
\definecolor{RurDynamixGreen}{RGB}{76, 175, 80}
\definecolor{RurDynamixBlue}{RGB}{33, 150, 243}
\definecolor{DarkGray}{RGB}{64, 64, 64}

\setbeamercolor{primary}{fg=white,bg=RurDynamixGreen}
\setbeamercolor{secondary}{fg=white,bg=RurDynamixBlue}
\setbeamertemplate{frametitle}{
  \begin{beamercolorbox}[rounded=true,shadow=true,leftskip=1cm,rightskip=1cm,wd=\paperwidth,ht=0.8cm,dp=0.4cm]{frametitle}
    \usebeamerfont{frametitle}\insertframetitle
  \end{beamercolorbox}
}

% Title Page Setup
\title{\tiny{Rural Territories in Transition: Morocco and the Global South}}
\author{Said Moukrim}
\institute{UM5-FSR Morocco}
\date{November 28, 2026}

\begin{document}

\begin{frame}

\begin{columns}[T]
  \column{0.5\textwidth}
  \includegraphics[width=\textwidth]{images/rur_dyn.png}

  \column{0.25\textwidth}
  \includegraphics[width=\textwidth]{images/FSR logo.png}
\end{columns}

\vspace{1cm}

\begin{center}
\Large {Hydraulic Innovation and Rural Transitions:}\\
\vspace{0.8cm}
\large{Spatio-Temporal Analysis of Artificial Basins Dynamics in Bab Berred}\\
Central Rif (Morocco)\\
\vspace{1.5cm}
\Large{Said Moukrim}
\end{center}

\end{frame}

% ==================== INTRODUCTION ====================

% Slide 2: Hydrological Context and Infrastructure Response
\begin{frame}
\frametitle{Hydrological Deficit and Infrastructure Adaptation}

\textbf{Hydroclimatic constraints in Mediterranean mountain regions:}
\begin{itemize}
  \item Increasing precipitation deficit in Rif Mountains
  \item Pronounced seasonal water stress during dry season
  \item Agricultural intensification requires reliable water supply
\end{itemize}

\vspace{0.3cm}

\textbf{Infrastructure-based adaptation strategy:}
\begin{itemize}
  \item Artificial basin proliferation as drought-mitigation infrastructure
  \item Water harvesting from surface runoff and groundwater resources
  \item Storage capacity expansion enabling irrigation during water-scarce periods
\end{itemize}

\vspace{0.3cm}

\textbf{Research question:}
\begin{itemize}
  \item Spatio-temporal dynamics and sustainability implications of basin expansion
  \item Linkages between infrastructure development, land-use change, and ecosystem response
\end{itemize}

\end{frame}

% Slide 3: Cannabis Cultivation Intensification
\begin{frame}
\frametitle{Agricultural Intensification: Hybrid Cannabis Cultivation}

\textbf{Market-driven agricultural transformation:}
\begin{itemize}
  \item Introduction of modern hybrid varieties (early 2010s) responding to international demand
  \item Productivity differential: traditional kif ($\approx$ 15 kg/ha) vs. hybrid varieties ($\approx$ 70 kg/ha)
  \item Four-to-five-fold yield increase explains rapid farmer adoption
\end{itemize}

\vspace{0.3cm}

\textbf{Hydrological consequence:}
\begin{itemize}
  \item Hybrid varieties require substantially greater water availability than traditional cultivation
  \item Necessitates transition from rainfed to irrigated agriculture
  \item Driver of artificial basin proliferation as adaptive infrastructure
\end{itemize}

\vspace{0.3cm}

\textbf{Note on legalization context:}
\begin{itemize}
  \item Law 13-21 (2021) created regulatory framework; not primary expansion driver
  \item International market demand, not legalization, shapes cultivation dynamics
\end{itemize}

\end{frame}

% Slide 4: Study Area and Regional Context
\begin{frame}
\frametitle{Study Site: Bab Berred Municipality, Chefchaouen Province}

\textbf{Geographic and climatic characteristics:}
\begin{itemize}
  \item Location: Central Rif Mountains, northern Morocco (35°00'N, 4°58'W)
  \item Area: 172.24 km² encompassing 17,224 hectares
  \item Topography: Steep terrain (mean slope 80\%), elevation 650-2,100 m
  \item Population: 24,742 residents dependent on agriculture
\end{itemize}

\vspace{0.3cm}

\textbf{Climate:}
\begin{itemize}
  \item Mediterranean mountain climate
  \item Mean annual precipitation $>$ 700 mm (winter concentration)
  \item Mean annual temperature: 18.6°C; June mean: 24.6°C
\end{itemize}

\vspace{0.3cm}

\begin{center}
  \vspace{1.2cm}
  \includegraphics[width=0.7\textwidth]{images/Figure1_Location.png}
  \small{\textit{Figure 1: Geographic location and topography of Bab Berred}}
\end{center}

\end{frame}

% ==================== METHODS ====================

% Slide 5: Methodology and Data Sources
\begin{frame}
\frametitle{Data Sources and Analytical Methods}

\textbf{Remote sensing data:}
\begin{itemize}
  \item Landsat 5, 7, 8 (30-m resolution, 1998-2025) via Google Earth Engine
  \item Historical imagery: Google Earth Pro, Esri World Imagery Wayback
  \item NDVI calculation from Landsat surface reflectance
\end{itemize}

\vspace{0.2cm}

\textbf{Climate data:}
\begin{itemize}
  \item Annual precipitation time series (1985-2025) from satellite products
  \item Trend analysis: Mann-Kendall test for monotonic trends
\end{itemize}

\vspace{0.2cm}

\textbf{Land-cover classification:}
\begin{itemize}
  \item ESA WorldCover 2021 product (seven classes identified in study area)
  \item Manual basin digitization and morphometric characterization
\end{itemize}

\vspace{0.2cm}

\textbf{Spatial analysis:}
\begin{itemize}
  \item QGIS for vector and raster processing
  \item Basin spatio-temporal analysis (2010-2025)
  \item Relationship analysis: basin distribution vs. topographic and land-use variables
\end{itemize}

\end{frame}

% ==================== RESULTS ====================

% Slide 6: Basin Proliferation Dynamics
\begin{frame}
\frametitle{Result 1: Artificial Basin Proliferation (2010-2025)}

\textbf{Basin inventory and growth trajectory:}
\begin{itemize}
  \item 2010: $\approx$ 28 basins
  \item 2025: 2,044 basins
  \item Overall basin density: 11.87 structures/km²
\end{itemize}

\vspace{0.3cm}

\textbf{Storage capacity expansion:}
\begin{itemize}
  \item 2010: $\approx$ 50,000 m³ aggregate volume
  \item 2025: $\approx$ 4.3 million m³ aggregate volume
  \item 85-fold increase: Massive hydrological transformation
\end{itemize}

\vspace{0.3cm}

\textbf{Morphometric miniaturization:}
\begin{itemize}
  \item Mean basin size 2010: 1,260 m²
  \item Mean basin size 2025: 812 m²
  \item 90\% of basins (2025): $<$ 1,500 m²
\end{itemize}

\vspace{0.3cm}

\begin{center}
  \vspace{0.8cm}
  \includegraphics[width=0.75\textwidth]{images/Figure2_BasinGrowth.png}
  \small{\textit{Figure 2: Cumulative basin count and temporal growth pattern (2010-2025)}}
\end{center}

\end{frame}

% Slide 7: Spatial Distribution Analysis
\begin{frame}
\frametitle{Result 2: Spatial Distribution and Environmental Controls}

\textbf{Topographic controls on basin location:}
\begin{itemize}
  \item Altitude preference: 800-1,500 m (optimal balance between water demand and source proximity)
  \item Slope distribution: 71\% on slopes 5-15\%; absent on slopes $>$ 25\% (engineering constraints)
\end{itemize}

\vspace{0.3cm}

\textbf{Land-use distribution:}
\begin{itemize}
  \item 78\% within cultivated agricultural areas
  \item 18\% in grasslands
  \item Only 3\% in natural vegetation (minimal ecosystem conversion direct to basins)
\end{itemize}

\vspace{0.3cm}

\textbf{Interpretation:}
\begin{itemize}
  \item Basin distribution reflects adaptation to biophysical constraints
  \item Lacks evidence of centralized planning; indicates individual farmer optimization
\end{itemize}

\vspace{0.3cm}

\begin{center}
  \vspace{0.5cm}
  \includegraphics[width=0.75\textwidth]{images/Figure3_SpatialDistribution.png}
  \small{\textit{Figure 3: Basin spatial distribution relative to altitude, slope, and land-use}}
\end{center}

\end{frame}

% Slide 8: Vegetation Index Evolution
\begin{frame}
\frametitle{Result 3: Spatio-Temporal Evolution of Vegetation Index}

\textbf{NDVI$_{\text{March-Jul}}$ temporal dynamics (1998-2025):}
\begin{itemize}
  \item 1998-2012: Gradual increase (Sen's slope 0.0025 per year)
  \item 2012-2025: Marked acceleration in vegetation index values
  \item Significant homogeneity break detected (2012): $p = 4.34 \times 10^{-6}$; tau = 0.501
\end{itemize}

\vspace{0.3cm}

\textbf{Seasonal pattern:}
\begin{itemize}
  \item Peak NDVI amplitudes coincide with cannabis cultivation growing season (May-August)
  \item Suggests irrigation-enabled vegetation productivity enhancement
\end{itemize}

\vspace{0.3cm}

\begin{center}
  \vspace{0.9cm}
  \includegraphics[width=0.75\textwidth]{images/Figure4_NDVI.png}
  \small{\textit{Figure 4: NDVI$_{\text{March-Jul}}$ time series (1998-2025) showing inflection point}}
\end{center}

\end{frame}

% Slide 9: Precipitation Trend Analysis
\begin{frame}
\frametitle{Result 4: Long-Term Precipitation Decline}

\textbf{Annual precipitation trends (1985-2024):}
\begin{itemize}
  \item Significant negative trend: Mann-Kendall tau = $-0.1376$, $p < 0.05$
  \item Progressive decline in annual precipitation over 40-year period
  \item High interannual variability: Wet years $>$ 850 mm; dry years $<$ 410 mm
\end{itemize}

\vspace{0.3cm}

\textbf{Extreme events:}
\begin{itemize}
  \item 2023: Record precipitation low ($\approx$ 300 mm)
  \item Intensified hydrological pressure on agriculture
\end{itemize}

\vspace{0.3cm}

\textbf{Critical finding:}
\begin{itemize}
  \item Divergence between declining precipitation and improved vegetation index
  \item Demonstrates infrastructure-based compensation for hydroclimatic deficit
\end{itemize}

\vspace{0.3cm}

\begin{center}
  \vspace{0.5cm}
  \includegraphics[width=0.75\textwidth]{images/Figure5_Precipitation.png}
  \small{\textit{Figure 5: Annual precipitation time series with Mann-Kendall trend}}
\end{center}

\end{frame}

% ==================== DISCUSSION ====================

% Slide 10: Vegetation Productivity Paradox and Ecological Reality
\begin{frame}
\frametitle{Discussion: Infrastructure-Induced Vegetation Improvement Paradox}

\textbf{Apparent vegetation productivity improvement masked ecological change:}
\begin{itemize}
  \item NDVI improvement indicates biomass productivity increase, not ecosystem health
  \item High NDVI signal ($>$ 0.75) derived from irrigated cannabis monoculture
  \item Irrigation infrastructure sustains vegetation productivity despite precipitation decline
\end{itemize}

\vspace{0.3cm}

\textbf{Contradictory land-use dynamics:}
\begin{itemize}
  \item Natural areas (grassland, rangeland) demonstrating regression
  \item Cultivated agricultural area expansion
  \item Ecosystem conversion, not ecosystem recovery
\end{itemize}

\vspace{0.3cm}

\textbf{Hydrological sustainability concern:}
\begin{itemize}
  \item Water extraction from basins and aquifer systems exceeds annual groundwater recharge capacity
  \item Aquifer overexploitation creating piezometric decline
  \item Long-term water security undermined despite short-term irrigation gains
\end{itemize}

\end{frame}

% Slide 11: Cumulative Ecosystem and Governance Implications
\begin{frame}
\frametitle{Discussion: Landscape Fragmentation and Governance Challenges}

\textbf{Ecosystem and biodiversity impacts:}
\begin{itemize}
  \item Landscape fragmentation from basin infrastructure and land-use conversion
  \item Disruption of wildlife habitat connectivity
  \item Endemic Mediterranean species persistence threatened
  \item Cumulative impacts of 2,000+ small reservoirs on local hydrology
\end{itemize}

\vspace{0.3cm}

\textbf{Adaptation paradox:}
\begin{itemize}
  \item Short-term: Infrastructure enables agricultural intensification and livelihood stability
  \item Long-term: Aquifer depletion dynamics create sustainability constraints
  \item Individual optimization (farmer basin construction) generates collective unsustainability
\end{itemize}

\vspace{0.3cm}

\textbf{Governance framework requirements:}
\begin{itemize}
  \item Sustainable extraction ceilings based on hydrological assessment
  \item Integrated territorial governance addressing water security, ecosystem integrity, livelihood resilience
  \item Infrastructure investment cannot substitute for comprehensive resource management
\end{itemize}

\end{frame}

% ==================== CONCLUSION ====================

% Slide 12: Conclusion and Global Implications
\begin{frame}
\frametitle{Conclusion: Rural Adaptation and Sustainability Challenges}

\textbf{Key findings:}
\begin{itemize}
  \item Spatio-temporal analysis reveals rapid artificial basin proliferation (28 $\to$ 2,044, 2010-2025)
  \item Infrastructure enables continued agricultural productivity despite declining precipitation
  \item Apparent vegetation improvement masks ecosystem transformation and aquifer overexploitation
\end{itemize}

\vspace{0.3cm}

\textbf{Regional and global relevance:}
\begin{itemize}
  \item Infrastructure-intensive agricultural transformation pattern replicated across Mediterranean and semi-arid regions
  \item Remote sensing-based environmental monitoring often misidentifies ecological status
  \item Integrated analysis (NDVI + precipitation + land-use + hydrology) essential for sustainability assessment
\end{itemize}

\vspace{0.3cm}

\textbf{Sustainability imperative:}
\begin{itemize}
  \item Sustainable rural futures require simultaneous attention to hydrological limits, ecosystem integrity, and livelihood resilience
  \item Particularly urgent as climate change intensifies water scarcity pressures
  \item Governance frameworks must transcend infrastructure investment alone
\end{itemize}

\end{frame}

% Final Slide
\begin{frame}[plain]
\begin{center}
  \vfill
  {\Huge \textbf{Thank you}}
  \\[1cm]
  {\Large Questions and Discussion}
  \\[2cm]
  \small{
    Said Moukrim, Ousmane Samba MAÏGA, et al. \\
    Mohammed-V University in Rabat \\
    Faculty of Sciences, Laboratory of Botany, Mycology and Environment
  }
  \vfill
\end{center}
\end{frame}

\end{document}
